% w07-tilt-geometry.tex
%   7주차 §7-1. 추진기 넷의 **자리**와 거치 부호 s_i, 그리고 틸팅 각 delta_i 가
%   각 유닛의 **자기 0도**에서 시작한다는 것.
%
%   축척 / scale
%       1 m = 32 mm. 평면도 — 화면 위가 선체 x_b, 화면 오른쪽이 y_b.
%       선체좌표 (x, y) 는 화면 (X, Y) = (32 y, 32 x) 에 놓인다.
%
%   수치의 출처 / where the numbers come from
%       `otter_config('quad_tilt')` 를 MATLAB 에서 그대로 찍은 것이다 (2026-11-03):
%           pos   = [ 1.5  1.5  -1.0  -1.0 ;  -0.395  0.395  -0.395  0.395 ]
%           mount = [ 1    1    -1    -1   ]        <- s_i
%           y_pont = 0.395,  l_f = 1.5,  l_a = 1.0,  alpha_max = 45 deg
%       눈대중으로 PPT 에서 베끼지 않는다.
%
%   거치 부호 / the mounting sign                (사용자 지시 2026-11-03)
%       "앞 뒤의 2개 모터가 바라보는 방향이 달라서 서보 각도의 정의가 달라."
%       선수 둘은 delta = 0 이 **+x_b**(앞)를 밀고, 선미 둘은 **-x_b**(뒤)를 민다.
%       그래서 s = [+1 +1 -1 -1] 이고, B 의 열이 그 부호를 들고 있다.
%       각도의 **부호 방향**은 둘이 같다 — 0 에서 y_b 쪽으로, 곧 이 평면도에서
%       시계방향이다. 다른 것은 **0 의 자리** 하나다. 자세한 것은
%       w07-servo-zero.tex 가 칸을 나눠 그린다.
%
%   회전 방향 / the sense of rotation            (CLAUDE.md 4-6-2)
%       z_b 가 아래이므로 +N 은 선수를 우현으로 돌리고, 위에서 내려다본 이 그림
%       (x_b 위, y_b 오른쪽)에서 그것은 **시계방향**이다. TikZ 의 arc 는 끝 각이
%       **작을 때** 시계이므로 전부 내려가는 각으로 적었다 —
%           unit 2:  arc (90:50:9)     0 도(화면 위)에서 우현 쪽으로
%           unit 4:  arc (270:230:9)   0 도(화면 아래)에서 우현 쪽으로
%           tau_N:   arc (330:210:7)
%       화면 기저로 풀어 오른손 법칙을 적용해 확인했다: 선수 위, 우현 오른쪽,
%       z_b 화면 안쪽 -> +N 은 위에서 오른쪽으로, 시계방향.
%
%   갑판을 네 자리가 **품도록** 그린 것 / the deck carries the stations
%       l_f = 1.5 m 는 otter.m 의 선체(L = 2.0 m, x in [-1, +1])보다 앞이다.
%       사용자 지시(2026-10-13)는 **유닛을 선체 밖에 그리지 말라**는 것이므로,
%       여기서는 폰툰을 네 자리가 들어가는 길이로 그리고 otter.m 의 선체 끝을
%       파선 하나로 밝힌다. 실습의 `W07_plot_scenario` 가 그리는 평면도와 **같은
%       처리**다 (xf = l_f + 0.25, xa = -l_a - 0.25). 6주차는 치수선만 뻗는
%       다른 처리를 했으나, 7주차는 유닛 **넷의 자리**가 그림의 주제라 자리를
%       그대로 보이는 쪽을 골랐다.

\documentclass[border=5pt]{standalone}
% gnc-style.tex — 이 볼트의 TikZ 그림이 공유하는 스타일.
%
% 저널 그림 규약을 스타일 하나로 굳혀 둔다. 그림마다 선 굵기나 화살촉을
% 다시 정하지 않는다 — 그것이 그림이 제각각으로 보이는 원인이다.
%
% 쓰는 법:  % gnc-style.tex — 이 볼트의 TikZ 그림이 공유하는 스타일.
%
% 저널 그림 규약을 스타일 하나로 굳혀 둔다. 그림마다 선 굵기나 화살촉을
% 다시 정하지 않는다 — 그것이 그림이 제각각으로 보이는 원인이다.
%
% 쓰는 법:  % gnc-style.tex — 이 볼트의 TikZ 그림이 공유하는 스타일.
%
% 저널 그림 규약을 스타일 하나로 굳혀 둔다. 그림마다 선 굵기나 화살촉을
% 다시 정하지 않는다 — 그것이 그림이 제각각으로 보이는 원인이다.
%
% 쓰는 법:  \input{gnc-style}  (figures/src/ 안에서)

\usepackage{tikz}
\usepackage{amsmath}
\usepackage{xcolor}
\usetikzlibrary{arrows.meta, positioning, calc, fit, backgrounds, decorations.pathmorphing}

% ---- 색 --------------------------------------------------------------------
% 색은 강조 하나에만 쓴다. 나머지는 검정.
\definecolor{ink}{HTML}{1A1A1A}
\definecolor{accent}{HTML}{7C3AED}
\definecolor{warn}{HTML}{B91C1C}
\definecolor{soft}{HTML}{666666}

% 기하 그림(좌표계·유도)에서만 쓰는 넷. 블록선도는 위의 흑백 + accent 로 충분하다.
\definecolor{pathc}{HTML}{0D9488}   % 계획 경로 · 항적
\definecolor{errc}{HTML}{C2410C}    % 오차 (cross-track 등)
\definecolor{alng}{HTML}{B45309}    % along-track
\definecolor{flow}{HTML}{0369A1}    % 조류 · 유속

% ---- 선과 화살촉 -----------------------------------------------------------
\tikzset{
  % 신호선. 화살촉은 작고 뾰족한 것 하나뿐이다.
  sig/.style   = {ink, line width=0.5pt, -{Stealth[length=4pt, width=3pt]}},
  wire/.style  = {ink, line width=0.5pt},
  % 강조 경로 — 파선으로 구분한다. 흑백 인쇄에서도 살아야 하기 때문이다.
  asig/.style  = {accent, line width=0.5pt, densely dashed,
                  -{Stealth[length=4pt, width=3pt]}},
  awire/.style = {accent, line width=0.5pt, densely dashed},
  % 블록. 흰 바탕, 직각, 검은 테두리.
  blk/.style   = {draw=ink, line width=0.55pt, fill=white,
                  minimum width=17mm, minimum height=8mm, inner sep=2pt, font=\small},
  % 합산점
  sum/.style   = {draw=ink, line width=0.55pt, fill=white, circle,
                  inner sep=0pt, minimum size=4.6mm, font=\scriptsize},
  asum/.style  = {draw=accent, line width=0.55pt, fill=white, circle,
                  inner sep=0pt, minimum size=4.6mm, font=\scriptsize},
  % 분기점
  dot/.style   = {ink, fill=ink, circle, inner sep=0pt, minimum size=1.5mm},
  % 신호 이름 — 선 위에 작게
  sname/.style = {font=\small, inner sep=1.5pt},
  % 그림 안의 짧은 설명. 그림 안의 글은 최소로 둔다
  note/.style  = {font=\scriptsize, soft, inner sep=1.5pt},
  anote/.style = {font=\scriptsize, accent, inner sep=1.5pt},
  % 축
  axis/.style  = {ink, line width=0.5pt},
  tick/.style  = {font=\scriptsize, soft},
}

% 캡션 한 줄 — 그림 아래 가는 선과 함께
\newcommand{\gnccaption}[2]{%
  \node[anchor=north west, text width=#1, align=left, font=\scriptsize, ink]
        at (current bounding box.south west) {\vspace{2pt}#2};}
  (figures/src/ 안에서)

\usepackage{tikz}
\usepackage{amsmath}
\usepackage{xcolor}
\usetikzlibrary{arrows.meta, positioning, calc, fit, backgrounds, decorations.pathmorphing}

% ---- 색 --------------------------------------------------------------------
% 색은 강조 하나에만 쓴다. 나머지는 검정.
\definecolor{ink}{HTML}{1A1A1A}
\definecolor{accent}{HTML}{7C3AED}
\definecolor{warn}{HTML}{B91C1C}
\definecolor{soft}{HTML}{666666}

% 기하 그림(좌표계·유도)에서만 쓰는 넷. 블록선도는 위의 흑백 + accent 로 충분하다.
\definecolor{pathc}{HTML}{0D9488}   % 계획 경로 · 항적
\definecolor{errc}{HTML}{C2410C}    % 오차 (cross-track 등)
\definecolor{alng}{HTML}{B45309}    % along-track
\definecolor{flow}{HTML}{0369A1}    % 조류 · 유속

% ---- 선과 화살촉 -----------------------------------------------------------
\tikzset{
  % 신호선. 화살촉은 작고 뾰족한 것 하나뿐이다.
  sig/.style   = {ink, line width=0.5pt, -{Stealth[length=4pt, width=3pt]}},
  wire/.style  = {ink, line width=0.5pt},
  % 강조 경로 — 파선으로 구분한다. 흑백 인쇄에서도 살아야 하기 때문이다.
  asig/.style  = {accent, line width=0.5pt, densely dashed,
                  -{Stealth[length=4pt, width=3pt]}},
  awire/.style = {accent, line width=0.5pt, densely dashed},
  % 블록. 흰 바탕, 직각, 검은 테두리.
  blk/.style   = {draw=ink, line width=0.55pt, fill=white,
                  minimum width=17mm, minimum height=8mm, inner sep=2pt, font=\small},
  % 합산점
  sum/.style   = {draw=ink, line width=0.55pt, fill=white, circle,
                  inner sep=0pt, minimum size=4.6mm, font=\scriptsize},
  asum/.style  = {draw=accent, line width=0.55pt, fill=white, circle,
                  inner sep=0pt, minimum size=4.6mm, font=\scriptsize},
  % 분기점
  dot/.style   = {ink, fill=ink, circle, inner sep=0pt, minimum size=1.5mm},
  % 신호 이름 — 선 위에 작게
  sname/.style = {font=\small, inner sep=1.5pt},
  % 그림 안의 짧은 설명. 그림 안의 글은 최소로 둔다
  note/.style  = {font=\scriptsize, soft, inner sep=1.5pt},
  anote/.style = {font=\scriptsize, accent, inner sep=1.5pt},
  % 축
  axis/.style  = {ink, line width=0.5pt},
  tick/.style  = {font=\scriptsize, soft},
}

% 캡션 한 줄 — 그림 아래 가는 선과 함께
\newcommand{\gnccaption}[2]{%
  \node[anchor=north west, text width=#1, align=left, font=\scriptsize, ink]
        at (current bounding box.south west) {\vspace{2pt}#2};}
  (figures/src/ 안에서)

\usepackage{tikz}
\usepackage{amsmath}
\usepackage{xcolor}
\usetikzlibrary{arrows.meta, positioning, calc, fit, backgrounds, decorations.pathmorphing}

% ---- 색 --------------------------------------------------------------------
% 색은 강조 하나에만 쓴다. 나머지는 검정.
\definecolor{ink}{HTML}{1A1A1A}
\definecolor{accent}{HTML}{7C3AED}
\definecolor{warn}{HTML}{B91C1C}
\definecolor{soft}{HTML}{666666}

% 기하 그림(좌표계·유도)에서만 쓰는 넷. 블록선도는 위의 흑백 + accent 로 충분하다.
\definecolor{pathc}{HTML}{0D9488}   % 계획 경로 · 항적
\definecolor{errc}{HTML}{C2410C}    % 오차 (cross-track 등)
\definecolor{alng}{HTML}{B45309}    % along-track
\definecolor{flow}{HTML}{0369A1}    % 조류 · 유속

% ---- 선과 화살촉 -----------------------------------------------------------
\tikzset{
  % 신호선. 화살촉은 작고 뾰족한 것 하나뿐이다.
  sig/.style   = {ink, line width=0.5pt, -{Stealth[length=4pt, width=3pt]}},
  wire/.style  = {ink, line width=0.5pt},
  % 강조 경로 — 파선으로 구분한다. 흑백 인쇄에서도 살아야 하기 때문이다.
  asig/.style  = {accent, line width=0.5pt, densely dashed,
                  -{Stealth[length=4pt, width=3pt]}},
  awire/.style = {accent, line width=0.5pt, densely dashed},
  % 블록. 흰 바탕, 직각, 검은 테두리.
  blk/.style   = {draw=ink, line width=0.55pt, fill=white,
                  minimum width=17mm, minimum height=8mm, inner sep=2pt, font=\small},
  % 합산점
  sum/.style   = {draw=ink, line width=0.55pt, fill=white, circle,
                  inner sep=0pt, minimum size=4.6mm, font=\scriptsize},
  asum/.style  = {draw=accent, line width=0.55pt, fill=white, circle,
                  inner sep=0pt, minimum size=4.6mm, font=\scriptsize},
  % 분기점
  dot/.style   = {ink, fill=ink, circle, inner sep=0pt, minimum size=1.5mm},
  % 신호 이름 — 선 위에 작게
  sname/.style = {font=\small, inner sep=1.5pt},
  % 그림 안의 짧은 설명. 그림 안의 글은 최소로 둔다
  note/.style  = {font=\scriptsize, soft, inner sep=1.5pt},
  anote/.style = {font=\scriptsize, accent, inner sep=1.5pt},
  % 축
  axis/.style  = {ink, line width=0.5pt},
  tick/.style  = {font=\scriptsize, soft},
}

% 캡션 한 줄 — 그림 아래 가는 선과 함께
\newcommand{\gnccaption}[2]{%
  \node[anchor=north west, text width=#1, align=left, font=\scriptsize, ink]
        at (current bounding box.south west) {\vspace{2pt}#2};}


\definecolor{hull}{HTML}{E6F0E6}
\definecolor{hulle}{HTML}{59735A}

%  폰툰 하나. #1 = 중심의 화면 X / one pontoon, centred on screen X = #1
\newcommand{\pontoon}[1]{%
  \filldraw[fill=hull, draw=hulle, line width=0.4pt]
    (#1-4.64,-40) -- (#1+4.64,-40) -- (#1+4.64,46.4) --
    (#1,56) -- (#1-4.64,46.4) -- cycle;}

%  선수 유닛. #1 = 화면 X / a bow unit: zero points up (+x_b)
\newcommand{\unitbow}[1]{%
  \begin{scope}[shift={(#1,48)}]
    \fill[accent, opacity=0.09] (0,0) -- (45:13) arc (45:135:13) -- cycle;
    \draw[accent, opacity=0.55, densely dashed, line width=0.35pt]
         (45:13) -- (0,0) -- (135:13);
    \draw[accent, line width=1pt, -{Stealth[length=5pt, width=4pt]}] (0,0) -- (0,15);
    \node[dot] at (0,0) {};
  \end{scope}}

%  선미 유닛. #1 = 화면 X / a stern unit: zero points down (-x_b)
\newcommand{\unitstern}[1]{%
  \begin{scope}[shift={(#1,-32)}]
    \fill[accent, opacity=0.09] (0,0) -- (225:13) arc (225:315:13) -- cycle;
    \draw[accent, opacity=0.55, densely dashed, line width=0.35pt]
         (225:13) -- (0,0) -- (315:13);
    \draw[accent, line width=1pt, -{Stealth[length=5pt, width=4pt]}] (0,0) -- (0,-15);
    \node[dot] at (0,0) {};
  \end{scope}}

%  치수선 하나 / one dimension line
\newcommand{\dimline}[4]{%
  \draw[soft, line width=0.4pt,
        {Stealth[length=3.5pt, width=2.5pt]}-{Stealth[length=3.5pt, width=2.5pt]}]
       (#1,#2) -- (#3,#4);}

\begin{document}
\begin{tikzpicture}[x=1mm, y=1mm, font=\scriptsize]

% ---------- 선체 / the hull -------------------------------------------------
\fill[hull, opacity=0.55] (-8,-40) rectangle (8,46.4);
\pontoon{-12.64}
\pontoon{12.64}
\draw[hulle, line width=0.4pt] (-12.64,0) -- (12.64,0);
\draw[soft, line width=0.3pt, densely dotted] (0,-48) -- (0,74);

%  otter.m 의 선체 끝. 선미 자리 x = -1.0 m 는 그 끝과 **같고**, 선수 자리
%  x = +1.5 m 만 앞으로 나간다 / only the forward station reaches past it
\draw[soft, densely dashed, line width=0.35pt] (-20,32) -- (20,32);
\node[note, anchor=east] at (-20.5,32) {otter.m hull ends, $x_b=+1$ m};

\node[note, anchor=south] at (0,58) {bow};
\node[note, anchor=north] at (0,-42) {stern};

% ---------- 네 유닛 / the four units ----------------------------------------
\unitbow{-12.64}
\unitbow{12.64}
\unitstern{-12.64}
\unitstern{12.64}

\node[ink, font=\scriptsize\bfseries, anchor=north, inner sep=1pt] at (-12.64,44.5) {1};
\node[ink, font=\scriptsize\bfseries, anchor=north, inner sep=1pt] at ( 12.64,44.5) {2};
\node[ink, font=\scriptsize\bfseries, anchor=south, inner sep=1pt] at (-12.64,-27.5) {3};
\node[ink, font=\scriptsize\bfseries, anchor=south, inner sep=1pt] at ( 12.64,-27.5) {4};

%  각의 부호 방향은 한 쌍에 하나씩만 그린다 — 넷을 다 그리면 읽히지 않는다
\draw[accent, line width=0.6pt, -{Stealth[length=4pt, width=3pt]}]
     (12.64,57) arc (90:50:9);
\node[anote, anchor=south west] at (17.2,57.8) {$\delta_2>0$};
\draw[accent, line width=0.6pt, -{Stealth[length=4pt, width=3pt]}]
     (12.64,-41) arc (270:230:9);
\node[anote, anchor=north west] at (17.2,-41.8) {$\delta_4>0$};

% ---------- 선체좌표계 / the body frame -------------------------------------
\draw[axis, -{Stealth[length=4pt, width=3pt]}] (0,0) -- (0,22);
\draw[axis, -{Stealth[length=4pt, width=3pt]}] (0,0) -- (22,0);
\node[ink, anchor=west]  at (1.6,20)  {$x_b$};
\node[ink, anchor=south] at (19,1.2)  {$y_b$};
\node[circle, draw=ink, fill=white, line width=0.4pt, inner sep=0.8pt] at (0,0) {};
\node[note, anchor=north east] at (-1.2,-1.2) {CO};

\begin{scope}[shift={(-28,-46)}]
  \draw[warn, line width=0.6pt, -{Stealth[length=4pt, width=3pt]}] (6.06,-3.5) arc (330:210:7);
  \node[warn, anchor=north, font=\scriptsize] at (0,-9) {$\tau_N>0$};
\end{scope}

% ---------- d_1 … d_4 -------------------------------------------------------
\draw[soft, line width=0.3pt, densely dotted] (-12.64,63.5) -- (-12.64,71);
\draw[soft, line width=0.3pt, densely dotted] ( 12.64,63.5) -- ( 12.64,71);
\dimline{-12.64}{70}{0}{70}
\dimline{0}{70}{12.64}{70}
\draw[soft, line width=0.4pt] (0,67.5) -- (0,72.5);
\node[note, anchor=south] at (-6.3,70.8) {$d_1$};
\node[note, anchor=south] at ( 6.3,70.8) {$d_2$};

\draw[soft, line width=0.3pt, densely dotted] (-12.64,-47.5) -- (-12.64,-55);
\draw[soft, line width=0.3pt, densely dotted] ( 12.64,-47.5) -- ( 12.64,-55);
\dimline{-12.64}{-54}{0}{-54}
\dimline{0}{-54}{12.64}{-54}
\draw[soft, line width=0.4pt] (0,-51.5) -- (0,-56.5);
\node[note, anchor=north] at (-6.3,-54.8) {$d_3$};
\node[note, anchor=north] at ( 6.3,-54.8) {$d_4$};

% ---------- l_f, l_a --------------------------------------------------------
\draw[soft, line width=0.3pt, densely dotted] (17.5,48) -- (29,48);
\draw[soft, line width=0.3pt, densely dotted] (0,0)    -- (29,0);
\draw[soft, line width=0.3pt, densely dotted] (17.5,-32) -- (29,-32);
\dimline{28}{0}{28}{48}
\dimline{28}{-32}{28}{0}
\draw[soft, line width=0.4pt] (26.5,48)  -- (29.5,48);
\draw[soft, line width=0.4pt] (26.5,0)   -- (29.5,0);
\draw[soft, line width=0.4pt] (26.5,-32) -- (29.5,-32);
\node[note, anchor=west] at (29.5,24)  {$l_f=1.50$ m};
\node[note, anchor=west] at (29.5,-16) {$l_a=1.00$ m};

% ---------- 자리 표 / the station table -------------------------------------
\node[ink, anchor=north west, align=left, font=\scriptsize] at (38,52) {%
  \begin{tabular}{@{}llrrr@{}}
    unit & station & $x_i$ [m] & $d_i=y_i$ [m] & $s_i$\\[1.5pt]
    1 & bow, port    & $+1.50$ & $-0.395$ & $+1$\\
    2 & bow, stbd    & $+1.50$ & $+0.395$ & $+1$\\
    3 & stern, port  & $-1.00$ & $-0.395$ & $-1$\\
    4 & stern, stbd  & $-1.00$ & $+0.395$ & $-1$\\
  \end{tabular}};

\node[note, anchor=north west, align=left, text width=46mm] at (38,22) {%
  $\delta_i$ runs from unit $i$'s own zero towards $y_b$, clockwise on this
  view, in both pairs.\\[2pt]
  $s_i=+1$: that zero points ahead.\quad $s_i=-1$: it points astern.\\[2pt]
  Shaded: the servo range $\lvert\delta_i\rvert\le 45^\circ$.};

\end{tikzpicture}
\end{document}
